\chapter{总结}

\note{生的 记得改 感觉没突出重点工作 减去前人算法的占比}


本文针对复杂环境下GNSS/INS松耦合组合导航系统中观测噪声呈现非高斯、重尾分布的问题，提出了一种融合自适应最大相关熵卡尔曼滤波（MCKF）的方法，并完成了完整的算法库开发与工程实现。主要工作总结如下：

\begin{enumerate}
\item \textbf{松耦合误差状态建模与滤波框架构建}：建立了15维INS误差状态动力学模型，包含位置误差、速度误差、失准角、陀螺零偏及加速度计零偏。采用误差状态卡尔曼滤波（ESKF）架构，将INS机械编排的名义状态与误差状态分离，实现了松耦合组合导航的基本滤波流程。该框架为后续MCKF的引入提供了清晰的模型基础。

\item \textbf{自适应MCKF与松耦合的融合方法}：将最大相关熵准则（MCC）引入松耦合卡尔曼滤波，替代传统的最小均方误差（MMSE）估计，利用相关熵在核空间中定义的局部相似性度量，有效利用了误差的高阶统计信息，增强了对非高斯、重尾观测噪声的鲁棒性。进一步提出了基于新息协方差匹配的自适应核带宽策略，分别对状态预测残差和观测残差动态调整核带宽，使滤波器能够在时变噪声环境下自适应地平衡预测与观测的置信度。

\item \textbf{模块化代码工程实现}：基于Python自主开发了可复现的组合导航算法库，采用“中强度接口”设计，实现了误差状态的分块管理与函数注入的模型扩展机制。项目采用事件驱动的双速率更新调度，支持IMU高频预测与GNSS稀疏校正的协同运行。误差传播采用矩阵指数离散化方法，陀螺与加计偏置以一阶马尔可夫过程建模，保证了数值一致性与参数可配置性。算法库已开源，支持不同滤波算法的横向对比与实测数据迁移。

\item \textbf{实验验证与后续展望}：通过仿真与实测数据验证了所提LC+自适应MCKF方法在观测异常、信号短暂中断等复杂场景下相比传统卡尔曼滤波具有更高的状态估计精度与系统鲁棒性。未来工作可进一步扩展至紧耦合架构，或融合视觉、激光雷达等多源传感器，结合深度学习提升全场景智能适应能力。
\end{enumerate}

本文的研究为复杂环境下高精度GNSS/INS组合导航提供了一种结构清晰、鲁棒性强的松耦合解决方案，所开发的代码工程为后续研究与实际部署提供了可复用的基础平台。